\begin{abstract}
An LLM asked to forget specific information should not lose knowledge or
capabilities that appear unrelated on the surface. Yet approximate
unlearning can cause such collateral damage. A single average score can
hide this damage when failures concentrate in a small but coherent subset
of retained behaviors. We therefore study whether susceptible retained
behaviors can be identified before unlearning and protected under a fixed
intervention budget. We model susceptibility using two complementary
signals: direction-agnostic local loss sensitivity and hidden-state
proximity to the forget set. For the gradient component, we define
\emph{loss-shake energy} as a pool-shared finite-difference estimator of
classical block-local gradient energy. Unlike
forget-update alignment, loss-shake energy is independent of the direction
ultimately taken by an unlearning trajectory. Paired perturbations estimate
it for an entire candidate pool using batched forward evaluations without
candidate-side backpropagation, while the representation component uses
forward-only hidden-state neighbors. We seal their joint profile before an
unchanged parent unlearner runs, admit only checkpoints passing an external
forgetting gate, and test two independent uses: held-out damage prediction and
fixed-budget repair allocation. Across multiple model families, objectives,
and benchmarks, the evaluation separately requires estimator fidelity,
positive prediction beyond both components and simple controls, lower mean
and worst-tail damage than four allocation controls, and non-inferior native
retained behavior under common forgetting and utility constraints. This
probe-to-protection framework links pre-unlearning risk interpretation to
targeted preservation without changing the parent objective or using realized
damage to select a checkpoint or repair pool.
\end{abstract}
